%%%%%%%%%%%%%%%%%%%%%%%%%%%%%%%%
% Chapter 5. Conclusion
%%%%%%%%%%%%%%%%%%%%%%%%%%%%%%%%

\section{Conclusion}\label{sec:conclusion}
\enlargethispage*{2\baselineskip}

We studied Random Soft Prompts (RSPs): training-free isotropic Gaussian vectors fitted to pretrained embedding statistics. The per-layer RSP contribution is controlled by an attention-mass scalar whose fixed-gap upper bound shrinks under KV cache growth, and isotropic resampling reaches top-$K$ entry regions that rank-preserving temperature cannot (\S\ref{sec:why_rsp_works}). Latent-prompt methods share the same entropy signature; the gain vanishes when one RSP is shared across samples (\S\ref{sec:experimental_analysis}); DAPO + RSP yields a modest late-training advantage (\S\ref{sec:application}). Because the signature emerges from purely random input-side perturbation, part of what trained \textit{soft prompts} deliver appears to be a \emph{structural by-product of injection}, and RSP composes orthogonally with sampling-based pipelines.
